\begin{problem}[Stalkwise isomorphisms are sheaf isomorphisms]
Let $\varphi:\mathcal F\to\mathcal G$ be a morphism of sheaves of abelian groups. Assume every $\varphi_x$ is an isomorphism. Construct the inverse sheaf morphism explicitly by local lifting and gluing.
\end{problem}

\paragraph{Why this problem matters.}
This is the strongest basic theorem showing that stalks completely detect local equivalence of sheaves.

\paragraph{Useful ideas before starting.}
Lift each germ of a target section, represent the lift on a neighborhood, shrink until it maps to the target section, then glue the local lifts.

\paragraph{Guided hints.}
On overlaps, use stalkwise injectivity first pointwise and then sheaf locality to prove the local lifts agree.

\ifdefined\FullProblemDossiers
\begin{solution}
Let $t\in\mathcal G(U)$. For each $x\in U$, choose $a_x\in\mathcal F_x$ with $\varphi_x(a_x)=t_x$. Represent $a_x$ by $s_x\in\mathcal F(V_x)$. Equality of germs lets us shrink $V_x$ so that $\varphi(s_x)=t|_{V_x}$. On $V_x\cap V_y$, the images of $s_x$ and $s_y$ agree. For each point of the overlap, stalkwise injectivity gives equality of their germs, so by the section-detection theorem they agree on the whole overlap. The $s_x$ therefore glue uniquely to $s\in\mathcal F(U)$ with $\varphi(s)=t$. The same injectivity argument shows this lift is unique. Define $\psi_U(t)=s$. Uniqueness makes the maps compatible with restriction, so $\psi$ is a sheaf morphism and is inverse to $\varphi$.
\end{solution}

\paragraph{What happened mathematically.}
Stalkwise algebra plus sheaf gluing reconstructs a global inverse morphism.

\paragraph{Extensions.}
Compare with AG2 Problem 7, which is deliberately reserved for VI/16 in its injective/surjective/exactness formulation.

\paragraph{Provenance.}
New dossier around a legacy theorem; AG2 Problem 7 remains reserved for VI/16 and is not duplicated as a numbered legacy problem here.
\else
\paragraph{Provenance.}
New dossier around a legacy theorem; AG2 Problem 7 remains reserved for VI/16 and is not duplicated as a numbered legacy problem here.
\fi
